\section{Introduction}
\label{sec:intro}

Humans come to know an unfamiliar physical world by acting in it. They push, drop, and collide objects, form an understanding of the principles behind what they see, and then carry that understanding back to vision: a glance at a new scene is enough to picture how it will unfold. We ask whether interaction gives a multimodal agent the same ability. The question is not whether the agent gets better at a task it practiced. An agent can improve over a few trials by searching a small action space, replaying a trajectory that worked, or exploiting a visual cue that happens to correlate with success. The question is whether what the agent learns by interacting with a world becomes an understanding of that world, one that still holds when the agent looks at a new scene from it.

Current benchmarks leave this distinction partly hidden. Passive physical-reasoning benchmarks ask models to judge physical concepts or outcomes from fixed images, videos, or questions; they do not test what feedback from action changes. Interactive physics benchmarks such as PHYRE~\cite{phyre} and I-PHYRE~\cite{iphyre}, and their recent VLM-agent counterparts~\cite{deepphy,ipr1}, measure task success, but success on the acted-on instance can mix physical inference with search. Indeed, VLMs trained through interaction improve on their training task without acquiring physical intuitions that transfer to related tasks~\cite{schulzebuschoff2026}. Scientific-discovery benchmarks~\cite{physgym,newtonbench,discoverphysics} let agents probe worlds with hidden or altered laws, but their observations and outputs are numeric or symbolic. Visual benchmarks of hidden physics~\cite{comphy,inphyre} test transfer to new scenes, but supply the evidence instead of letting the agent collect it. What remains open is the full path from interaction to vision: whether an agent that has experimented in a physical world can foresee how a new scene of that world will look.

We introduce a benchmark that tests this path. Each \emph{latent physical world} is indexed by four factors: a physical-phenomenon family, a world-parameter regime, a latent-dynamics condition, and a latent-law realization. Each realization contains a set of counterfactual prediction instances that share its laws and constants. In an episode, the agent first receives a fixed budget of tool-mediated interaction on support instances. Tool calls expose observable state, such as object attributes, positions, velocities, environmental measurements, and time, but never reveal the hidden constants, governing equation, latent transformation, or answer key. After support, feedback is frozen. For each held-out instance of the same realization, the agent must select the correct future frame from four rendered candidates.

This design separates an understanding of the world from success on the scenes the agent practiced on. Support and held-out instances share the same laws and constants but differ in initial state, target object, and prediction interval, so a support trajectory cannot be replayed. Candidate positions are independently permuted, and distractors are rendered from nearby but physically distinct successor states. Two controls make the transfer argument testable rather than assumed. A \emph{cross-realization} control supplies support evidence from a different hidden law of the same kind; if accuracy does not drop, the support evidence was not used as evidence about the particular law. A \emph{canonical-physics} baseline predicts with the untransformed family law and identifies instances that can be answered without inferring the latent transformation.

The benchmark contains 5,670 instances: seven physical-phenomenon families (central-force, kinematics, fluid, contact, dissipative, field-interaction, and elastic dynamics), three world-parameter regimes per family, three latent-dynamics conditions per regime, six latent-law realizations per condition, and fifteen instances per realization. The protocol is organized around three questions. Does budgeted support improve held-out prediction (RQ1)? Is the improvement specific to the latent law, and does it hold across structural, coefficient-field, and combined latent changes (RQ2)? Where does the chain break: perception, evidence gathering, hypothesis formation, or visual decoding (RQ3)? For RQ3, a separate measurement probe asks the agent to estimate a hidden intrinsic attribute, such as mass or friction, with the basic experiment tools, and scores both the accuracy of the estimate and the number of tool calls it took.

Our contributions are:
\begin{itemize}
  \item A controlled benchmark for testing whether agent-selected physical experiments support held-out counterfactual visual prediction under an unchanged latent law.
  \item A 5,670-instance implementation spanning seven physical-phenomenon families \TODO{Rescope to the families validated end to end; non-central families currently ignore the hidden law in their ground truth.}, with simulator-generated ground truth, rendered candidates, and per-instance candidate permutation.
  \item A stratified evaluation protocol with built-in controls (support-budget sweep, same- versus cross-realization evaluation, appearance and canonical-physics baselines, and law-informed ceilings) and a measurement probe that scores attribute-estimation accuracy and tool-call cost, applied to frontier multimodal agents at the model--harness level. \TODO{Replace this clause with the main empirical finding once results land.}
\end{itemize}

\begin{figure*}[t]
  \centering
  % \includegraphics[width=\textwidth]{figures/fig1_hero.pdf}
  \fbox{\parbox[c][6cm][c]{0.95\textwidth}{\centering
    \textbf{Figure 1 placeholder (hero figure).}\\[0.5em]
    \emph{Support:} the agent spends a fixed tool budget on several instances from one latent-law realization and records observable state.\\
    \emph{Evaluation:} tool access is frozen; the agent sees a new initial state from the same realization and chooses among four rendered future frames.\\
    \emph{Contrast:} appearance matching to the support history fails when the held-out state changes; an understanding of the world's laws can transfer.
  }}
  \caption{\textbf{Benchmark protocol.} The benchmark separates interaction on support instances from counterfactual prediction on held-out states. During support, the agent gathers observable evidence about one latent-law realization. During evaluation, tool access is frozen and the agent must choose the correct future frame for a new state governed by the same laws and constants.}
  \label{fig:hero}
\end{figure*}

Section~\ref{sec:related} positions the benchmark relative to prior work. Section~\ref{sec:benchmark} formalizes the task. Section~\ref{sec:construction} describes the benchmark construction, tool interface, simulator, and splits. Section~\ref{sec:protocol} defines the evaluation protocol and controls. Section~\ref{sec:validation} gives validation checks, Section~\ref{sec:experiments} reports baselines and analyses, and Section~\ref{sec:conclusion} discusses limitations and conclusions.
